\setcounter{secnumdepth}{-1}

\chapter{Introduction}

Aujourd’hui plus que jamais, les discussions sur les enjeux de la biodiversité se multiplient. Toutefois, la mobilisation n’est pas nouvelle : dès le début des années 1990, la communauté internationale s’est engagée en faveur de la protection des espèces, avec l’adoption de la Convention des Nations Unies sur la diversité biologique (Secretariat of the Convention on Biological Diversity, 2025), dont le but est de préserver la biodiversité, qui est au cœur des débats actuels. 

Cette perte de biodiversité se traduit par un déclin préoccupant de nombreuses espèces. Selon l'édition 2025 de la liste rouge établie par l’\gls{uicn}, sur les 172 620 espèces évaluées, plus de 48 600, soit 28\%, sont à l’heure actuelle menacées d’extinction (The IUCN Red List of Threatened Species, 2025). 

Dans l’imaginaire collectif européen, l’érosion de la biodiversité est généralement associée à la disparition des grands mammifères exotiques tels que les tigres, lions, rhinocéros, etc. Mais, qu’en est-il des espèces plus proches de nos régions ? 

En Europe, au moins 1 677 espèces, soit 11\% des espèces évaluées par l’\gls{uicn}, sont menacées d’extinction (The IUCN Red List of Threatened Species, 2025).

Notre projet s’est intéressé au hérisson européen (Erinaceus europaeus), une espèce largement appréciée du public mais encore insuffisamment étudiée et documentée, notamment en ce qui concerne l’état de sa population. E. europaeus est un petit mammifère nocturne insectivore de la famille des érinacéidés (Erinaceidae). Reconnaissable à ses piquants, le hérisson est reconnu pour son rôle d’espèce parapluie, c’est-à-dire que sa protection permet de maintenir de nombreuses autres espèces animales (amphibiens, insectes, petits mammifères) ainsi que la flore, qui partagent le même habitat (Roberge et al., 2004). En plus de ce rôle, il contribue à la régulation des invertébrés et à l’équilibre de son écosystème. 

Le hérisson voit sa population reculer dans plusieurs régions du continent européen (Gazzard et al., 2025). Les causes sont multiples. Parmi celles-ci, on peut citer : fragmentation croissante des habitats, urbanisation, intensification agricole, usage massif de pesticides ou encore mortalité routière (figure 1). Ces pressions cumulées fragilisent davantage cette espèce caractérisée par une hibernation sensible aux perturbations, la saisonnalité de sa reproduction, une faible capacité de dispersion, ainsi qu’une sensibilité accrue face à certaines maladies. Autant de traits qui limitent sa résilience face aux changements environnementaux. Cette situation a récemment conduit l’\gls{uicn} (The IUCN Red List of Threatened Species, 2025) à revoir le classement du hérisson européen, qui passe de la catégorie “préoccupation mineure” à la catégorie des espèces “quasi menacées”, reflétant une inquiétude croissante quant à son avenir à l’échelle continentale (The IUCN Red List of Threatened Species, 2012). 

Dans le cadre de notre projet, nous nous sommes intéressés à la situation en Belgique. En Flandre, les données issues d’observations citoyennes indiquent une chute de près de 50\% du nombre d’individus recensés entre 2010 et 2018 (Gazzard et al., 2024). En Wallonie, il existe peu de données ou d’études scientifiques qui permettent de dresser un bilan démographique. Faute de données structurées et de suivis systématiques, l’état réel de la population demeure difficile à évaluer (figure Ce 2). Ce déficit d’informations constitue aujourd’hui un frein majeur à l’élaboration de stratégies de conservation cohérentes et adaptées dans cette région de Belgique. Il semble donc indispensable de combler les lacunes en matière de données. Dans un premier temps, il faut comprendre les raisons de ces lacunes, puis, dans un second temps, proposer des solutions pour résoudre le problème.  

Pour les besoins de notre projet, nous avons passé en revue les données disponibles en Flandre, au Luxembourg, dans certaines régions de France ainsi qu’en Wallonie afin d’établir une base de connaissances permettant de comparer la situation dans ces pays et régions avec celle de la Wallonie. Nous avons évalué la manière dont les acteurs concernés (pouvoirs publics, organisations non gouvernementales et populations locales) procèdent au recensement du hérisson ainsi que le fonctionnement des différents outils utilisés à cette fin. L’objectif final est de mettre en évidence les solutions qui sont efficaces dans les régions et pays étudiés et de déterminer si des solutions similaires pourraient être appliquées en Wallonie. Parallèlement à la question du recensement, nous avons rassemblé des informations présentant la biologie générale du hérisson (caractéristiques physiques, habitat, prédateurs, nutrition, reproduction, etc.). Ces informations nous paraissant indispensables pour mieux comprendre cette espèce et donc favoriser sa protection. 

Nous avons ensuite effectué des recherches pour déterminer les raisons du déclin de la population du hérisson, les conséquences de son déclin sur le plan environnemental et socio-économique ainsi que les solutions envisageables pour pallier ces problèmes.  

Toutes les recherches détaillées que nous avons effectuées se fondent sur des articles scientifiques, des sites naturalistes, ainsi que sur des rapports gouvernementaux. Nous avons également consulté des blogs et des forums pour évaluer la connaissance générale et le degré de conscience du grand public sur le hérisson. Pour nous aider dans nos recherches, nous avons également exploré des applications de recensement déjà existantes.  

Ces recherches nous ont menés à créer un dispositif de collecte d’observations géoréférencées innovant, qui viendra compléter le travail d’inventaire des professionnels. 